\label{sec:abstract}
Large Language Models (LLMs) fundamentally reshape spam detection by enabling both synthetic training data generation and AI-powered spam attacks. We systematically evaluate this dual impact through comprehensive experimental tracks covering synthetic-to-real training scenarios, zero-shot detection of AI-generated attacks, and cross-model augmentation strategies. Our rigorous evaluation across multiple LLMs, prompt variations, and classifier architectures reveals three critical findings. First, synthetic spam can partially substitute for real training data with dramatic model-dependent quality gaps: high-quality generation maintains 86\% of baseline performance while low-quality generation retains only 54\%. Second, zero-shot detection of AI-generated spam achieves moderate baseline performance, with substantial variation across generation sources. Third, cross-model synthetic augmentation significantly improves detection effectiveness, with diversity-driven training providing larger gains than quality-optimized approaches. Ensemble classifiers consistently demonstrate superior robustness to distributional shifts compared to linear models. These findings reveal asymmetric advantages favoring attackers through accessible generation capabilities, though defenders can leverage robust architectures and cross-model augmentation to maintain detection effectiveness. Our work provides empirical foundations for understanding Generative AI's role in cybersecurity offense-defense dynamics.
